% =====================================================================
%  sec4.tex — 4. Ramification
%  原文 pp.5–7
% =====================================================================
\section{分歧}

不过，在回到叠之前，我想先说明一下几何 Langlands 语境中所谓的「分歧」是什么意思。

对我们迄今所说的内容，一个简单的推广是：考虑的不是闭的定向二维流形 $C$ 上的平坦丛，而是带穿刺点的二维流形 $C'=C\setminus p$ 上的平坦丛；也就是说，$C'$ 是从 $C$ 中挖去一点 $p$ 而得到的\zhnt{4.1}。

我们选取一个共轭类 $\mathcal{C}\subset G_{\C}$，并记 $\YC(G,C';\mathcal{C})$ 为同态 $\rho:\pi_{1}(C')\to G_{\C}$ 的模空间（模去共轭），要求绕 $p$ 的单值性落在共轭类 $\mathcal{C}$ 中\zhnt{4.2}。

我们以前作过的许多陈述，在这种带穿刺点的情形下都有自然的类比。特别地，$\YC(G,C';\mathcal{C})$ 具有复辛流形的自然结构：它有自然的复结构与全纯辛形式 $\Omega$。与无穿刺点的情形一样，我们可以定义 $\YC(G,C';\mathcal{C})$ 的 B-模型；同时，把 $\YC(G,C';\mathcal{C})$ 看作以 $\omega=\operatorname{Im}\Omega$ 为辛形式的实辛流形，我们也可以定义 A-模型。B-模型与 A-模型都完全不依赖 $C'$ 的复结构。

其次，引入对偶群 $\LG$，并令 ${}^{L}\mathcal{C}$ 为其复化中的一个共轭类。同样，空间 $\YC(\LG,C';{}^{L}\mathcal{C})$ 有自然的 B-模型与 A-模型。

基于我们迄今所说的，人们或许会猜想：对 $\mathcal{C}$ 与 ${}^{L}\mathcal{C}$ 之间的某个映射，$\YC(G,C';\mathcal{C})$ 与 $\YC(\LG,C';{}^{L}\mathcal{C})$ 之间可能存在镜像对称。这个问题的答案是「不完全是」，原因有几条。一个问题是：$G_{\C}$ 中的共轭类与 ${}^{L}G_{\C}$ 中的共轭类之间没有自然的对应。更根本的问题是：$\YC(G,C';\mathcal{C})$ 的 B-模型全纯地依赖于共轭类 $\mathcal{C}$，而同一空间的 A-模型却不\zhnt{4.3}。要找到一种有可能正确的陈述形式，我们必须添加额外的参数，才能得到一个镜像对称的参数组。

无论如何，不管添加什么参数，只要把 $C'$ 仅看成一个带穿刺点的定向二维流形，镜像对称的问题就极难回答。为了取得进展，正如在非分歧情形（即无穿刺点的情形）中那样，给 $C'$ 配上一个复结构并使用 Hitchin 方程是非常有帮助的。这实际上也帮助我们找到正确的参数：在试图给带穿刺点曲面上的 Hitchin 方程一个自然表述时，一套改进了的参数就出现了。设 $z$ 是穿刺点附近的局部参数，写作 $z=re^{i\theta}$。在带穿刺点的情形，自然地（见 Simpson 1990）引入取值于极大环面 $T\subset G$ 的李代数 $\mathfrak{t}$ 的变量 $\alpha,\beta,\gamma$，并考虑 $C'$ 上 Hitchin 方程的解，使其在 $z=0$ 附近的行为如下：
\begin{equation}
A=\alpha\,d\theta+\cdots
\label{eq:A-asym}
\end{equation}
\begin{equation}
\phi=\beta\,\frac{dr}{r}-\gamma\,d\theta+\cdots.
\label{eq:phi-asym}
\end{equation}
省略号指在 $z=0$ 附近奇性更弱的项\zhnt{4.4}。

关于 Hitchin 方程的一切通常陈述，在这种情形下都有密切的类比。以这种奇性为奇性的 Hitchin 方程之解的模空间，是一个超凯勒流形 $\MH(G,C;\alpha,\beta,\gamma)$。在其中一个复结构（通常记作 $J$）下，它与 $\YC(G,C;\mathcal{C})$ 一致，其中 $\mathcal{C}$ 是包含 $U=\exp(-2\pi(\alpha-i\gamma))$ 的共轭类\yz{为简单起见，我们假定 $U$ 是正则的。一般情形下成立的、更为复杂的陈述见 Gukov 与 Witten (2006)。}\zhnt{4.5}。在另一个复结构（常记作 $I$）下，$\MH(G,C;\alpha,\beta,\gamma)$ 参数化 Higgs 丛 $(E,\varphi)$，其中 $\varphi\in H^{0}(C',K\otimes\ad(E))$ 在 $z=0$ 处有极点，且
\[
\varphi\sim\tfrac12(\beta+i\gamma)\,(dz/z).
\]
此外还有 Hitchin 纤维化，并且非分歧情形的大部分通常陈述——既包括我们已解释过的，也包括此处省略未讲的——都有密切的类比。更详细的解释以及对原始文献的指引，见 Gukov 与 Witten (2006)。

变量 $\alpha,\beta,\gamma$ 是经典几何的一套自然参数。然而在量子层面，还有一个自然变量 $\eta$（类似于规范场论中通常的 $\theta$-角），如 Gukov 与 Witten (2006) 第 2.3 节所述。有了完整的参数组 $(\alpha,\beta,\gamma,\eta)$，就有可能表述一个自然的对偶陈述：$\MH(\LG,C;{}^{L}\alpha,{}^{L}\beta,{}^{L}\gamma,{}^{L}\eta)$ 镜像于 $\MH(G,C;\alpha,\beta,\gamma,\eta)$，参数之间有某个映射，见 Gukov 与 Witten (2006) 第 2.4 节。这个映射的要点是 $(\alpha,\eta)=({}^{L}\eta,-{}^{L}\alpha)$\zhnt{4.6}。由于单值性 $U$ 依赖于 ${}^{L}\alpha$，这就说明：单值性的对偶涉及量子参数 $\eta$，而后者在经典几何中是看不见的。在 A-模型中，$\eta$ 成为复化 Kähler 类的虚部。

这一对偶陈述在再次把 A-膜对应到 D-模之后，给出这一情形下几何 Langlands 对偶的一个陈述，与 Frenkel 与 Gaitsgory (2005) 中通过代数几何与二维共形场论所得的结果类似。

\begin{remark}
我们在此停下，解释关于经典几何的一个很基本的事实，它可作为第 5 节的背景。在复结构 $J$ 下，具有式 \eqref{eq:A-asym} 那种奇性的 Hitchin 方程之解，描述一个平坦 $G_{\C}$-丛 $E\to C'$，其绕穿刺点 $p$ 有单值性。$E$ 可以延拓过 $p$ 成为一个全纯丛，当然不能延拓为平坦丛；而且，从全纯的观点看，$E$ 在 $p$ 附近可以被平凡化。于是 $E\to C'$ 上的平坦联络在这一规范下，由 $C'$ 上的一个全纯 $(1,0)$-形式（取值于 $G_{\C}$ 的李代数）表示，该形式在 $p$ 处有单极点：
\begin{equation}
A=dz\left(\frac{\alpha-i\gamma}{iz}+\cdots\right),
\label{eq:regular}
\end{equation}
其中省略的项在 $z=0$ 处正则。联络在 $z=0$ 处的奇性之所以是单极点，是因为 Hitchin 方程的拟设 \eqref{eq:A-asym} 只允许 $1/|z|$ 阶的奇性。具有这种单极点的全纯联络，称为有正则奇点\zhnt{4.7}。

在几何 Langlands 中，通常所谓的驯顺分歧，粗略地说就是我们刚刚到达的情形：全纯丛 $E\to C$ 带有一个有正则奇点的全纯联络形式。其实，「驯顺分歧」这个说法有时被用来指单极点的留数是幂零的情形，而在式 \eqref{eq:regular} 中我们似乎处在相反的情形，即半单留数。Gukov 与 Witten (2006) 中说明：只要稍加小心，$\MH(G,C';\alpha,\beta,\gamma,\eta)$ 的镜像对称实际上就是一个足够好的框架，可以理解任意留数的单极点情形下的几何 Langlands。例如，幂零留数的情形可以通过令 ${}^{L}\alpha={}^{L}\gamma=0$（或在对偶描述中令 $\gamma=\eta=0$）来理解。
\end{remark}
